\documentclass[11pt,a4paper]{article}
\usepackage[margin=2.2cm]{geometry}
\usepackage{booktabs,amsmath,hyperref,enumitem,xcolor}
\hypersetup{colorlinks,urlcolor=blue!50!black,linkcolor=blue!50!black}
\setlist{itemsep=1pt,topsep=3pt}
\newcommand{\cost}[1]{\hfill\textcolor{gray}{\small[#1]}}
\title{\textbf{CrackDistill --- What to do next, and why}\\\large An honest diagnosis and ranked ideas}
\author{Shahin Alakbarov}\date{28 October 2026}
\begin{document}\maketitle

\section{Diagnosis: why KD does not help (yet)}
Knowledge distillation helps when the teacher knows something the student cannot learn from the labels.
Our measurements say SAM\,2 does not:
\begin{itemize}
\item \textbf{The teacher is worse than the labels.} Zero-shot SAM\,2 against the ground truth: median IoU \textbf{0.605}. Recall is 0.56 on cracks 4\,px wide or thinner, and 21\% of its crack pixels are ground-truth background.
\item \textbf{The student has already caught up.} The student alone reaches a test pixel F1 of 0.71--0.73. If the teacher is not clearly better than that, it only adds noise or label smoothing.
\item \textbf{Augmentation does the same job.} The v1 KD gain (+0.053) came while the student had no spatial augmentation. Flips and rotations alone gave +0.104, and KD added 0 on top of them.
\item \textbf{The teacher cannot follow the strongest augmentation.} Cached logits cannot follow mosaic or scale changes, and that is exactly where the best model (A4, 0.540) gets its advantage.
\end{itemize}
\textbf{Consequence:} a SAM\,1+SAM\,2 fusion of two \emph{zero-shot} teachers is unlikely to fix this. Their errors are
correlated: both miss thin cracks, and both are prompted with the same boxes. Fusion is worth trying only after the teacher has been improved.

\paragraph{This is not failure.} ``A carefully controlled study shows that foundation-model KD gains on a
small crack segmenter vanish under matched augmentation'' is a publishable negative result. Many KD papers compare
against weak baselines, as our own v1 did. The ideas below either turn this into a positive result or make the
negative result stronger.

\section{Ranked ideas}

\subsection*{1. Make the teacher better than the student (highest expected gain)\cost{$\sim$4 T4-h}}
Fine-tune SAM\,2 on the Crack500 train crops, either the normalisation layers only (Segment Any Crack, arXiv:2504.14138)
or with LoRA/adapters (CrackSAM, arXiv:2312.04233). Re-cache the logits and rerun A5.
\textbf{Gate before any KD run:} the teacher's val pixel F1 must clearly exceed the student's (above 0.80, against the student's 0.73).
If the teacher passes, KD has real knowledge to transfer. If not, we have learned that cheaply.

\subsection*{2. Use SAM where there are no labels: semi-supervised scaling\cost{$\sim$6 T4-h}}
The teacher's value is in \emph{labelling new images}, not in softening labels we already have.
\begin{itemize}
\item \textbf{Unlabelled pool:} the 250 full train photos at native scale (including crack-free regions), and images from other
pavement sets with their labels set aside.
\item \textbf{Pseudo-labels:} the fine-tuned teacher (idea 1), keeping only confident pixels. Then train A4 on labelled plus pseudo-labelled
images, mean-teacher or UniMatch style (SemiSAM, S4M, Boxes2Pixels).
\item \textbf{Crack-free tiles:} these also fix our main full-photo error, false positives on bare road (precision about 0.69).
\end{itemize}

\subsection*{3. Test where a foundation teacher should matter: out-of-distribution\cost{$\sim$1 h, no training}}
We have only measured in-domain. A SAM teacher may help \emph{generalisation} even with no in-domain gain.
Evaluate the existing A1, A4 and A5 checkpoints zero-shot on DeepCrack (already local), CFD and GAPs384, using pixel F1 and clDice.
\textbf{Cheapest experiment on this list, and it could change the story.}

\subsection*{4. Distil from a teacher that can follow full augmentation\cost{$\sim$6 T4-h}}
Train a large YOLOv11x-seg with A4's augmentation (optionally on SAM pseudo-labels), then distil it into YOLOv11n online,
in the same batch as the student, so mosaic and scale are no problem.
A chain such as SAM $\rightarrow$ YOLO-x $\rightarrow$ YOLO-n (a ``teacher assistant'') closes the capacity gap and the augmentation gap at once.

\subsection*{5. Fix the student's architecture for thin cracks\cost{$\sim$3 T4-h}}
YOLO-seg predicts masks from prototypes at 1/4 resolution (128$\times$128 for a 512 input), so a 3\,px crack is under one prototype cell.
Options:
\begin{itemize}
\item a P2 (stride-4) head or a 2$\times$ prototype resolution;
\item \texttt{mask\_ratio}=2;
\item a training input of 640.
\end{itemize}
This improves the baseline whether or not KD works. Compare against a lightweight semantic crack model (SCSegamba, CVPR 2025).

\subsection*{6. Better data protocol\cost{$\sim$2 T4-h}}
\begin{itemize}
\item \textbf{Combined training set.} Crack500 train + DeepCrack train + CFD half, with clean splits, then test on held-out sets. More diverse data usually beats any loss trick.
\item \textbf{Train on full photos, not only crops.} The 250 full train photos are in an older Kaggle mirror. Use random native-scale tiles, including empty ones.
\item \textbf{Use the right metric.} Pixel F1 and clDice instead of instance mask mAP50, which over-penalises fragmented cracks (test mAP 0.27 against pixel F1 0.71).
\end{itemize}

\subsection*{7. SAM\,1 + SAM\,2 fusion, done properly (after idea 1)\cost{$\sim$4 T4-h}}
\begin{itemize}
\item Standardise each teacher's logits (Sun et al., CVPR 2024), then weight them per pixel by confidence (low entropy).
\item Apply KD only where the teachers agree with each other and with the ground truth; elsewhere use the ground truth.
\item Use the fusion mainly to produce \emph{pseudo-labels} for idea 2, where an ensemble really reduces noise.
\item Schedule: KD weight $w(e)=w_0\cos^2(\pi e/2E)$, strong early and fading late. Our logs show KD helps only as an early regulariser.
\end{itemize}

\section{Recommended plan (2--3 weeks)}
\begin{center}\small
\begin{tabular}{clll}
\toprule
Week & Do & Decides & If it fails\\
\midrule
1 & Idea 3 (OOD eval) + E2 rerun + seeds 1--2 & Is there any KD signal? & Strengthens the negative result\\
1 & Idea 1 (fine-tune the teacher), check teacher F1 & Can the teacher beat the student? & Skip KD, go to 5/6\\
2 & Idea 2 or 4 with the good teacher & Positive KD result vs A4 & Report the negative result\\
2--3 & Idea 5 + 6 (architecture, data) & Best final model & --\\
\bottomrule
\end{tabular}\end{center}
\textbf{Success criterion, fixed in advance:} beat A4 by at least 0.02 val mask mAP50 (or 0.02 pixel F1) with the same sign on test,
confirmed over 3 seeds.

\paragraph{The paper either way.}
\begin{itemize}
\item \textbf{Positive:} ``A crack-tuned SAM teacher with semi-supervised distillation improves a 2.8M-parameter
segmenter beyond strong augmentation.''
\item \textbf{Negative:} ``Zero-shot foundation-model KD gains are an artefact of weak baselines'', with the audit, controls and data-efficiency curve.
\end{itemize}
Both stories rest on the work already done.
\end{document}
